% Options for packages loaded elsewhere
\PassOptionsToPackage{unicode}{hyperref}
\PassOptionsToPackage{hyphens}{url}
\documentclass[
]{article}
\usepackage{xcolor}
\usepackage[margin=1in]{geometry}
\usepackage{amsmath,amssymb}
\setcounter{secnumdepth}{-\maxdimen} % remove section numbering
\usepackage{iftex}
\ifPDFTeX
  \usepackage[T1]{fontenc}
  \usepackage[utf8]{inputenc}
  \usepackage{textcomp} % provide euro and other symbols
\else % if luatex or xetex
  \usepackage{unicode-math} % this also loads fontspec
  \defaultfontfeatures{Scale=MatchLowercase}
  \defaultfontfeatures[\rmfamily]{Ligatures=TeX,Scale=1}
\fi
\usepackage{lmodern}
\ifPDFTeX\else
  % xetex/luatex font selection
\fi
% Use upquote if available, for straight quotes in verbatim environments
\IfFileExists{upquote.sty}{\usepackage{upquote}}{}
\IfFileExists{microtype.sty}{% use microtype if available
  \usepackage[]{microtype}
  \UseMicrotypeSet[protrusion]{basicmath} % disable protrusion for tt fonts
}{}
\makeatletter
\@ifundefined{KOMAClassName}{% if non-KOMA class
  \IfFileExists{parskip.sty}{%
    \usepackage{parskip}
  }{% else
    \setlength{\parindent}{0pt}
    \setlength{\parskip}{6pt plus 2pt minus 1pt}}
}{% if KOMA class
  \KOMAoptions{parskip=half}}
\makeatother
\usepackage{color}
\usepackage{fancyvrb}
\newcommand{\VerbBar}{|}
\newcommand{\VERB}{\Verb[commandchars=\\\{\}]}
\DefineVerbatimEnvironment{Highlighting}{Verbatim}{commandchars=\\\{\}}
% Add ',fontsize=\small' for more characters per line
\usepackage{framed}
\definecolor{shadecolor}{RGB}{248,248,248}
\newenvironment{Shaded}{\begin{snugshade}}{\end{snugshade}}
\newcommand{\AlertTok}[1]{\textcolor[rgb]{0.94,0.16,0.16}{#1}}
\newcommand{\AnnotationTok}[1]{\textcolor[rgb]{0.56,0.35,0.01}{\textbf{\textit{#1}}}}
\newcommand{\AttributeTok}[1]{\textcolor[rgb]{0.13,0.29,0.53}{#1}}
\newcommand{\BaseNTok}[1]{\textcolor[rgb]{0.00,0.00,0.81}{#1}}
\newcommand{\BuiltInTok}[1]{#1}
\newcommand{\CharTok}[1]{\textcolor[rgb]{0.31,0.60,0.02}{#1}}
\newcommand{\CommentTok}[1]{\textcolor[rgb]{0.56,0.35,0.01}{\textit{#1}}}
\newcommand{\CommentVarTok}[1]{\textcolor[rgb]{0.56,0.35,0.01}{\textbf{\textit{#1}}}}
\newcommand{\ConstantTok}[1]{\textcolor[rgb]{0.56,0.35,0.01}{#1}}
\newcommand{\ControlFlowTok}[1]{\textcolor[rgb]{0.13,0.29,0.53}{\textbf{#1}}}
\newcommand{\DataTypeTok}[1]{\textcolor[rgb]{0.13,0.29,0.53}{#1}}
\newcommand{\DecValTok}[1]{\textcolor[rgb]{0.00,0.00,0.81}{#1}}
\newcommand{\DocumentationTok}[1]{\textcolor[rgb]{0.56,0.35,0.01}{\textbf{\textit{#1}}}}
\newcommand{\ErrorTok}[1]{\textcolor[rgb]{0.64,0.00,0.00}{\textbf{#1}}}
\newcommand{\ExtensionTok}[1]{#1}
\newcommand{\FloatTok}[1]{\textcolor[rgb]{0.00,0.00,0.81}{#1}}
\newcommand{\FunctionTok}[1]{\textcolor[rgb]{0.13,0.29,0.53}{\textbf{#1}}}
\newcommand{\ImportTok}[1]{#1}
\newcommand{\InformationTok}[1]{\textcolor[rgb]{0.56,0.35,0.01}{\textbf{\textit{#1}}}}
\newcommand{\KeywordTok}[1]{\textcolor[rgb]{0.13,0.29,0.53}{\textbf{#1}}}
\newcommand{\NormalTok}[1]{#1}
\newcommand{\OperatorTok}[1]{\textcolor[rgb]{0.81,0.36,0.00}{\textbf{#1}}}
\newcommand{\OtherTok}[1]{\textcolor[rgb]{0.56,0.35,0.01}{#1}}
\newcommand{\PreprocessorTok}[1]{\textcolor[rgb]{0.56,0.35,0.01}{\textit{#1}}}
\newcommand{\RegionMarkerTok}[1]{#1}
\newcommand{\SpecialCharTok}[1]{\textcolor[rgb]{0.81,0.36,0.00}{\textbf{#1}}}
\newcommand{\SpecialStringTok}[1]{\textcolor[rgb]{0.31,0.60,0.02}{#1}}
\newcommand{\StringTok}[1]{\textcolor[rgb]{0.31,0.60,0.02}{#1}}
\newcommand{\VariableTok}[1]{\textcolor[rgb]{0.00,0.00,0.00}{#1}}
\newcommand{\VerbatimStringTok}[1]{\textcolor[rgb]{0.31,0.60,0.02}{#1}}
\newcommand{\WarningTok}[1]{\textcolor[rgb]{0.56,0.35,0.01}{\textbf{\textit{#1}}}}
\usepackage{longtable,booktabs,array}
\usepackage{calc} % for calculating minipage widths
% Correct order of tables after \paragraph or \subparagraph
\usepackage{etoolbox}
\makeatletter
\patchcmd\longtable{\par}{\if@noskipsec\mbox{}\fi\par}{}{}
\makeatother
% Allow footnotes in longtable head/foot
\IfFileExists{footnotehyper.sty}{\usepackage{footnotehyper}}{\usepackage{footnote}}
\makesavenoteenv{longtable}
\usepackage{graphicx}
\makeatletter
\newsavebox\pandoc@box
\newcommand*\pandocbounded[1]{% scales image to fit in text height/width
  \sbox\pandoc@box{#1}%
  \Gscale@div\@tempa{\textheight}{\dimexpr\ht\pandoc@box+\dp\pandoc@box\relax}%
  \Gscale@div\@tempb{\linewidth}{\wd\pandoc@box}%
  \ifdim\@tempb\p@<\@tempa\p@\let\@tempa\@tempb\fi% select the smaller of both
  \ifdim\@tempa\p@<\p@\scalebox{\@tempa}{\usebox\pandoc@box}%
  \else\usebox{\pandoc@box}%
  \fi%
}
% Set default figure placement to htbp
\def\fps@figure{htbp}
\makeatother
\setlength{\emergencystretch}{3em} % prevent overfull lines
\providecommand{\tightlist}{%
  \setlength{\itemsep}{0pt}\setlength{\parskip}{0pt}}
\usepackage{bookmark}
\IfFileExists{xurl.sty}{\usepackage{xurl}}{} % add URL line breaks if available
\urlstyle{same}
\hypersetup{
  pdftitle={Chapter 4: Distribution Fitting},
  hidelinks,
  pdfcreator={LaTeX via pandoc}}

\title{Chapter 4: Distribution Fitting}
\author{}
\date{\vspace{-2.5em}}

\begin{document}
\maketitle

\begin{Shaded}
\begin{Highlighting}[]
\NormalTok{data }\OtherTok{\textless{}{-}} \FunctionTok{read.csv}\NormalTok{(}\StringTok{"../02\_data\_preparation/int\_data.csv"}\NormalTok{)}
\end{Highlighting}
\end{Shaded}

\begin{Shaded}
\begin{Highlighting}[]
\NormalTok{dates }\OtherTok{\textless{}{-}} \FunctionTok{as.Date}\NormalTok{(data}\SpecialCharTok{$}\NormalTok{date, }\AttributeTok{format =} \StringTok{"\%Y{-}\%m{-}\%d"}\NormalTok{)}
\NormalTok{alior }\OtherTok{\textless{}{-}}\NormalTok{ data}\SpecialCharTok{$}\NormalTok{closes\_alior}
\NormalTok{ing }\OtherTok{\textless{}{-}}\NormalTok{ data}\SpecialCharTok{$}\NormalTok{closes\_ing}
\end{Highlighting}
\end{Shaded}

\subsection{Fitting a Distribution to
Data}\label{fitting-a-distribution-to-data}

One of the fundamental steps in data analysis is identifying a
theoretical distribution that best describes the empirical data. This
enables the subsequent application of statistical inference methods,
tests, or predictive models that assume a specific distributional
structure. In this chapter, several popular continuous distributions
were fitted to data characterized by a right-skewed tail.

\subsubsection{What is distribution
fitting?}\label{what-is-distribution-fitting}

Distribution fitting is the process of finding a theoretical
distribution (e.g., normal, log-normal, gamma, Weibull) that best
describes empirical data.\\
The goals are: - estimating the distribution parameters based on the
sample,\\
- assessing the quality of the fit,\\
- selecting the best model for further analysis and inference.

\subsubsection{Selected distributions}\label{selected-distributions}

\paragraph{a) Log-normal distribution}\label{a-log-normal-distribution}

If a random variable \(X\) follows a normal distribution:

\[
X \sim N(\mu, \sigma^2),
\]

then the variable \(Y = e^X\) follows a log-normal distribution. The
probability density function (pdf) is:

\[
f(y; \mu, \sigma) = \frac{1}{y \, \sigma \sqrt{2\pi}} \exp\left(-\frac{(\ln y - \mu)^2}{2\sigma^2}\right), \quad y > 0.
\]

\textbf{Parameters:} - \(\mu\) -- mean of the logarithms,\\
- \(\sigma\) -- standard deviation of the logarithms.

\paragraph{b) Gamma distribution}\label{b-gamma-distribution}

Probability density function:

\[
f(x; k, \theta) = \frac{1}{\Gamma(k)\theta^k} x^{k-1} e^{-x/\theta}, \quad x > 0.
\]

\textbf{Parameters:} - \(k\) -- shape parameter,\\
- \(\theta\) -- scale parameter.\\
\(\Gamma(k)\) is the gamma function.

\paragraph{c) Weibull distribution}\label{c-weibull-distribution}

Probability density function:

\[
f(x; k, \lambda) = \frac{k}{\lambda} \left(\frac{x}{\lambda}\right)^{k-1} e^{-(x/\lambda)^k}, \quad x > 0.
\]

\textbf{Parameters:} - \(k\) -- shape parameter,\\
- \(\lambda\) -- scale parameter.

\subsubsection{Parameter estimation --
MLE}\label{parameter-estimation-mle}

\textbf{Maximum Likelihood Estimator (MLE)} selects such distribution
parameters that maximize the likelihood function:

\[
L(\theta | x_1, \dots, x_n) = \prod_{i=1}^n f(x_i; \theta),
\]

where \(\theta\) is the vector of distribution parameters.

Often, the log-likelihood function is used:

\[
\ell(\theta) = \ln L(\theta | x_1, \dots, x_n) = \sum_{i=1}^n \ln f(x_i; \theta).
\]

MLE finds such \(\hat{\theta}\) for which \(\ell(\theta)\) reaches its
maximum.

\begin{center}\rule{0.5\linewidth}{0.5pt}\end{center}

\subsubsection{\texorpdfstring{\texttt{fitdistrplus}
package}{fitdistrplus package}}\label{fitdistrplus-package}

The \textbf{fitdistrplus} package allows: - parameter estimation for
many distributions (e.g., using MLE),\\
- visualization of the fit (histograms, QQ-plots, CDF),\\
- goodness-of-fit testing,\\
- model comparison using information criteria (AIC, BIC).

\begin{Shaded}
\begin{Highlighting}[]
\CommentTok{\# install.packages("fitdistrplus")}
\FunctionTok{library}\NormalTok{(fitdistrplus)}
\end{Highlighting}
\end{Shaded}

\begin{verbatim}
## Loading required package: MASS
\end{verbatim}

\begin{verbatim}
## Loading required package: survival
\end{verbatim}

\subsubsection{Data Preparation}\label{data-preparation}

The analyzed data were transformed so that all values are positive. To
achieve this, the entire dataset was shifted by adding the value
\(|min(\text{x})| + 0.001\), ensuring the absence of negative values and
allowing the use of distributions defined only on the positive domain
(e.g., Gamma, Weibull, Log-normal).

\begin{Shaded}
\begin{Highlighting}[]
\NormalTok{alior\_rev }\OtherTok{\textless{}{-}} \SpecialCharTok{{-}}\NormalTok{alior }\SpecialCharTok{+} \FunctionTok{abs}\NormalTok{(}\FunctionTok{min}\NormalTok{(}\SpecialCharTok{{-}}\NormalTok{alior)) }\SpecialCharTok{+} \FloatTok{0.001}
\end{Highlighting}
\end{Shaded}

\begin{Shaded}
\begin{Highlighting}[]
\NormalTok{ing\_rev }\OtherTok{\textless{}{-}} \SpecialCharTok{{-}}\NormalTok{ing }\SpecialCharTok{+} \FunctionTok{abs}\NormalTok{(}\FunctionTok{min}\NormalTok{(}\SpecialCharTok{{-}}\NormalTok{ing)) }\SpecialCharTok{+} \FloatTok{0.001}
\end{Highlighting}
\end{Shaded}

\subsubsection{Methods}\label{methods}

Distribution fitting was performed using Maximum Likelihood Estimation
(MLE). Parameters were estimated using the fitdistr function. Three
distributions were fitted:

\begin{itemize}
\tightlist
\item
  log-normal,
\item
  gamma,
\item
  Weibull.
\end{itemize}

For each distribution, the parameters were estimated and theoretical
curves were fitted to the histogram of the empirical data. The quality
of the fit was then evaluated both visually and using statistical tests
and information criteria.

\begin{Shaded}
\begin{Highlighting}[]
\NormalTok{alior\_fit\_lnorm\_rev }\OtherTok{\textless{}{-}} \FunctionTok{fitdist}\NormalTok{(alior\_rev, }\StringTok{"lnorm"}\NormalTok{)}
\FunctionTok{summary}\NormalTok{(alior\_fit\_lnorm\_rev)}
\end{Highlighting}
\end{Shaded}

\begin{verbatim}
## Fitting of the distribution ' lnorm ' by maximum likelihood 
## Parameters : 
##         estimate Std. Error
## meanlog 2.905137 0.03957731
## sdlog   1.070052 0.02798527
## Loglikelihood:  -3210.394   AIC:  6424.787   BIC:  6433.976 
## Correlation matrix:
##         meanlog sdlog
## meanlog       1     0
## sdlog         0     1
\end{verbatim}

\begin{Shaded}
\begin{Highlighting}[]
\NormalTok{pdf\_alior\_lnorm }\OtherTok{\textless{}{-}} \ControlFlowTok{function}\NormalTok{(x) \{}
  \FunctionTok{dlnorm}\NormalTok{(}
    \AttributeTok{x =} \FunctionTok{abs}\NormalTok{(}\FunctionTok{min}\NormalTok{(}\SpecialCharTok{{-}}\NormalTok{alior)) }\SpecialCharTok{+} \FloatTok{0.001} \SpecialCharTok{{-}}\NormalTok{ x,}
    \AttributeTok{meanlog =}\NormalTok{ alior\_fit\_lnorm\_rev}\SpecialCharTok{$}\NormalTok{estimate[}\StringTok{"meanlog"}\NormalTok{],}
    \AttributeTok{sdlog =}\NormalTok{ alior\_fit\_lnorm\_rev}\SpecialCharTok{$}\NormalTok{estimate[}\StringTok{"sdlog"}\NormalTok{]}
\NormalTok{  )}
\NormalTok{\}}
\end{Highlighting}
\end{Shaded}

\begin{Shaded}
\begin{Highlighting}[]
\NormalTok{alior\_fit\_gamma\_rev }\OtherTok{\textless{}{-}} \FunctionTok{fitdist}\NormalTok{(alior\_rev, }\StringTok{"gamma"}\NormalTok{)}
\FunctionTok{summary}\NormalTok{(alior\_fit\_gamma\_rev)}
\end{Highlighting}
\end{Shaded}

\begin{verbatim}
## Fitting of the distribution ' gamma ' by maximum likelihood 
## Parameters : 
##         estimate  Std. Error
## shape 1.33826251 0.063126911
## rate  0.04824041 0.002747261
## Loglikelihood:  -3142.514   AIC:  6289.029   BIC:  6298.218 
## Correlation matrix:
##           shape      rate
## shape 1.0000000 0.8277001
## rate  0.8277001 1.0000000
\end{verbatim}

\begin{Shaded}
\begin{Highlighting}[]
\NormalTok{pdf\_alior\_gamma }\OtherTok{\textless{}{-}} \ControlFlowTok{function}\NormalTok{(x) \{}
  \FunctionTok{dgamma}\NormalTok{(}
    \AttributeTok{x =} \FunctionTok{abs}\NormalTok{(}\FunctionTok{min}\NormalTok{(}\SpecialCharTok{{-}}\NormalTok{alior)) }\SpecialCharTok{+} \FloatTok{0.001} \SpecialCharTok{{-}}\NormalTok{ x,}
    \AttributeTok{shape =}\NormalTok{ alior\_fit\_gamma\_rev}\SpecialCharTok{$}\NormalTok{estimate[}\StringTok{"shape"}\NormalTok{],}
    \AttributeTok{rate  =}\NormalTok{ alior\_fit\_gamma\_rev}\SpecialCharTok{$}\NormalTok{estimate[}\StringTok{"rate"}\NormalTok{]}
\NormalTok{  )}
\NormalTok{\}}
\end{Highlighting}
\end{Shaded}

\begin{Shaded}
\begin{Highlighting}[]
\NormalTok{alior\_fit\_weibull\_rev }\OtherTok{\textless{}{-}} \FunctionTok{fitdist}\NormalTok{(alior\_rev, }\StringTok{"weibull"}\NormalTok{)}
\FunctionTok{summary}\NormalTok{(alior\_fit\_weibull\_rev)}
\end{Highlighting}
\end{Shaded}

\begin{verbatim}
## Fitting of the distribution ' weibull ' by maximum likelihood 
## Parameters : 
##        estimate Std. Error
## shape  1.229412 0.03708809
## scale 29.634912 0.93886728
## Loglikelihood:  -3138.554   AIC:  6281.107   BIC:  6290.296 
## Correlation matrix:
##           shape     scale
## shape 1.0000000 0.3136991
## scale 0.3136991 1.0000000
\end{verbatim}

\begin{Shaded}
\begin{Highlighting}[]
\NormalTok{pdf\_alior\_weibull }\OtherTok{\textless{}{-}} \ControlFlowTok{function}\NormalTok{(x) \{}
  \FunctionTok{dweibull}\NormalTok{(}
    \AttributeTok{x =} \FunctionTok{abs}\NormalTok{(}\FunctionTok{min}\NormalTok{(}\SpecialCharTok{{-}}\NormalTok{alior)) }\SpecialCharTok{+} \FloatTok{0.001} \SpecialCharTok{{-}}\NormalTok{ x,}
    \AttributeTok{shape =}\NormalTok{ alior\_fit\_weibull\_rev}\SpecialCharTok{$}\NormalTok{estimate[}\StringTok{"shape"}\NormalTok{],}
    \AttributeTok{scale =}\NormalTok{ alior\_fit\_weibull\_rev}\SpecialCharTok{$}\NormalTok{estimate[}\StringTok{"scale"}\NormalTok{]}
\NormalTok{  )}
\NormalTok{\}}
\end{Highlighting}
\end{Shaded}

\begin{Shaded}
\begin{Highlighting}[]
\NormalTok{x }\OtherTok{\textless{}{-}} \FunctionTok{seq}\NormalTok{(}
  \AttributeTok{from =} \FunctionTok{min}\NormalTok{(alior),}
  \AttributeTok{to =} \FunctionTok{max}\NormalTok{(alior),}
  \AttributeTok{length.out =} \DecValTok{1000}
\NormalTok{)}

\FunctionTok{plot}\NormalTok{(}
  \AttributeTok{x =}\NormalTok{ x,}
  \AttributeTok{y =} \FunctionTok{pdf\_alior\_lnorm}\NormalTok{(x),}
  \AttributeTok{col =} \StringTok{"red"}\NormalTok{,}
  \AttributeTok{type =} \StringTok{"l"}\NormalTok{,}
  \AttributeTok{lwd =} \DecValTok{2}\NormalTok{,}
  \AttributeTok{ylim =} \FunctionTok{c}\NormalTok{(}\DecValTok{0}\NormalTok{, }\FunctionTok{max}\NormalTok{(}
    \FunctionTok{pdf\_alior\_lnorm}\NormalTok{(x),}
    \FunctionTok{pdf\_alior\_gamma}\NormalTok{(x),}
    \FunctionTok{pdf\_alior\_weibull}\NormalTok{(x)}
\NormalTok{  )),}
  \AttributeTok{xlab =} \StringTok{"x"}\NormalTok{,}
  \AttributeTok{ylab =} \StringTok{"Density"}\NormalTok{,}
  \AttributeTok{main =} \StringTok{"Estimated Theoretical Distributions for Alior Bank"}
\NormalTok{)}

\FunctionTok{lines}\NormalTok{(}
  \AttributeTok{x =}\NormalTok{ x,}
  \AttributeTok{y =} \FunctionTok{pdf\_alior\_gamma}\NormalTok{(x),}
  \AttributeTok{col =} \StringTok{"blue"}\NormalTok{,}
  \AttributeTok{lwd =} \DecValTok{2}
\NormalTok{)}

\FunctionTok{lines}\NormalTok{(}
  \AttributeTok{x =}\NormalTok{ x,}
  \AttributeTok{y =} \FunctionTok{pdf\_alior\_weibull}\NormalTok{(x),}
  \AttributeTok{col =} \StringTok{"darkgreen"}\NormalTok{,}
  \AttributeTok{lwd =} \DecValTok{3}
\NormalTok{)}

\FunctionTok{legend}\NormalTok{(}
  \StringTok{"topleft"}\NormalTok{,}
  \AttributeTok{legend=}\FunctionTok{c}\NormalTok{(}\StringTok{"Log{-}normal"}\NormalTok{, }\StringTok{"Gamma"}\NormalTok{, }\StringTok{"Weibull"}\NormalTok{),}
  \AttributeTok{col=}\FunctionTok{c}\NormalTok{(}\StringTok{"red"}\NormalTok{, }\StringTok{"blue"}\NormalTok{, }\StringTok{"darkgreen"}\NormalTok{),}
  \AttributeTok{lwd=}\DecValTok{2}
\NormalTok{)}
\end{Highlighting}
\end{Shaded}

\pandocbounded{\includegraphics[keepaspectratio]{04_distribution_fitting_files/figure-latex/unnamed-chunk-12-1.pdf}}

\begin{Shaded}
\begin{Highlighting}[]
\NormalTok{ing\_fit\_lnorm\_rev }\OtherTok{\textless{}{-}} \FunctionTok{fitdist}\NormalTok{(ing\_rev, }\StringTok{"lnorm"}\NormalTok{)}
\FunctionTok{summary}\NormalTok{(ing\_fit\_lnorm\_rev)}
\end{Highlighting}
\end{Shaded}

\begin{verbatim}
## Fitting of the distribution ' lnorm ' by maximum likelihood 
## Parameters : 
##          estimate Std. Error
## meanlog 4.0843974 0.03569409
## sdlog   0.9650614 0.02523941
## Loglikelihood:  -3996.942   AIC:  7997.883   BIC:  8007.072 
## Correlation matrix:
##         meanlog sdlog
## meanlog       1     0
## sdlog         0     1
\end{verbatim}

\begin{Shaded}
\begin{Highlighting}[]
\NormalTok{pdf\_ing\_lnorm }\OtherTok{\textless{}{-}} \ControlFlowTok{function}\NormalTok{(x) \{}
  \FunctionTok{dlnorm}\NormalTok{(}
    \AttributeTok{x =} \FunctionTok{abs}\NormalTok{(}\FunctionTok{min}\NormalTok{(}\SpecialCharTok{{-}}\NormalTok{ing)) }\SpecialCharTok{+} \FloatTok{0.001} \SpecialCharTok{{-}}\NormalTok{ x,}
    \AttributeTok{meanlog =}\NormalTok{ ing\_fit\_lnorm\_rev}\SpecialCharTok{$}\NormalTok{estimate[}\StringTok{"meanlog"}\NormalTok{],}
    \AttributeTok{sdlog   =}\NormalTok{ ing\_fit\_lnorm\_rev}\SpecialCharTok{$}\NormalTok{estimate[}\StringTok{"sdlog"}\NormalTok{]}
\NormalTok{  )}
\NormalTok{\}}
\end{Highlighting}
\end{Shaded}

\begin{Shaded}
\begin{Highlighting}[]
\NormalTok{ing\_fit\_gamma\_rev }\OtherTok{\textless{}{-}} \FunctionTok{fitdist}\NormalTok{(ing\_rev, }\StringTok{"gamma"}\NormalTok{)}
\FunctionTok{summary}\NormalTok{(ing\_fit\_gamma\_rev)}
\end{Highlighting}
\end{Shaded}

\begin{verbatim}
## Fitting of the distribution ' gamma ' by maximum likelihood 
## Parameters : 
##         estimate  Std. Error
## shape 1.74501515 0.083647410
## rate  0.02148073 0.001188585
## Loglikelihood:  -3887.802   AIC:  7779.604   BIC:  7788.792 
## Correlation matrix:
##           shape      rate
## shape 1.0000000 0.8631708
## rate  0.8631708 1.0000000
\end{verbatim}

\begin{Shaded}
\begin{Highlighting}[]
\NormalTok{pdf\_ing\_gamma }\OtherTok{\textless{}{-}} \ControlFlowTok{function}\NormalTok{(x) \{}
  \FunctionTok{dgamma}\NormalTok{(}
    \AttributeTok{x =} \FunctionTok{abs}\NormalTok{(}\FunctionTok{min}\NormalTok{(}\SpecialCharTok{{-}}\NormalTok{ing)) }\SpecialCharTok{+} \FloatTok{0.001} \SpecialCharTok{{-}}\NormalTok{ x,}
    \AttributeTok{shape =}\NormalTok{ ing\_fit\_gamma\_rev}\SpecialCharTok{$}\NormalTok{estimate[}\StringTok{"shape"}\NormalTok{],}
    \AttributeTok{rate  =}\NormalTok{ ing\_fit\_gamma\_rev}\SpecialCharTok{$}\NormalTok{estimate[}\StringTok{"rate"}\NormalTok{]}
\NormalTok{  )}
\NormalTok{\}}
\end{Highlighting}
\end{Shaded}

\begin{Shaded}
\begin{Highlighting}[]
\NormalTok{ing\_fit\_weibull\_rev }\OtherTok{\textless{}{-}} \FunctionTok{fitdist}\NormalTok{(ing\_rev, }\StringTok{"weibull"}\NormalTok{)}
\FunctionTok{summary}\NormalTok{(ing\_fit\_weibull\_rev)}
\end{Highlighting}
\end{Shaded}

\begin{verbatim}
## Fitting of the distribution ' weibull ' by maximum likelihood 
## Parameters : 
##        estimate Std. Error
## shape  1.490118 0.04554992
## scale 89.594147 2.33489855
## Loglikelihood:  -3872.894   AIC:  7749.787   BIC:  7758.976 
## Correlation matrix:
##           shape     scale
## shape 1.0000000 0.3053983
## scale 0.3053983 1.0000000
\end{verbatim}

\begin{Shaded}
\begin{Highlighting}[]
\NormalTok{pdf\_ing\_weibull }\OtherTok{\textless{}{-}} \ControlFlowTok{function}\NormalTok{(x) \{}
  \FunctionTok{dweibull}\NormalTok{(}
    \AttributeTok{x =} \FunctionTok{abs}\NormalTok{(}\FunctionTok{min}\NormalTok{(}\SpecialCharTok{{-}}\NormalTok{ing)) }\SpecialCharTok{+} \FloatTok{0.001} \SpecialCharTok{{-}}\NormalTok{ x,}
    \AttributeTok{shape =}\NormalTok{ ing\_fit\_weibull\_rev}\SpecialCharTok{$}\NormalTok{estimate[}\StringTok{"shape"}\NormalTok{],}
    \AttributeTok{scale =}\NormalTok{ ing\_fit\_weibull\_rev}\SpecialCharTok{$}\NormalTok{estimate[}\StringTok{"scale"}\NormalTok{]}
\NormalTok{  )}
\NormalTok{\}}
\end{Highlighting}
\end{Shaded}

\begin{Shaded}
\begin{Highlighting}[]
\NormalTok{x }\OtherTok{\textless{}{-}} \FunctionTok{seq}\NormalTok{(}
  \AttributeTok{from =} \FunctionTok{min}\NormalTok{(ing),}
  \AttributeTok{to =} \FunctionTok{max}\NormalTok{(ing),}
  \AttributeTok{length.out =} \DecValTok{1000}
\NormalTok{)}

\FunctionTok{plot}\NormalTok{(}
  \AttributeTok{x =}\NormalTok{ x,}
  \AttributeTok{y =} \FunctionTok{pdf\_ing\_lnorm}\NormalTok{(x),}
  \AttributeTok{col =} \StringTok{"red"}\NormalTok{,}
  \AttributeTok{type =} \StringTok{"l"}\NormalTok{,}
  \AttributeTok{lwd =} \DecValTok{2}\NormalTok{,}
  \AttributeTok{ylim =} \FunctionTok{c}\NormalTok{(}\DecValTok{0}\NormalTok{, }\FunctionTok{max}\NormalTok{(}
    \FunctionTok{pdf\_ing\_lnorm}\NormalTok{(x),}
    \FunctionTok{pdf\_ing\_gamma}\NormalTok{(x),}
    \FunctionTok{pdf\_ing\_weibull}\NormalTok{(x)}
\NormalTok{  )),}
  \AttributeTok{xlab =} \StringTok{"x"}\NormalTok{,}
  \AttributeTok{ylab =} \StringTok{"Density"}\NormalTok{,}
  \AttributeTok{main =} \StringTok{"Estimated Theoretical Distributions for ING Bank"}
\NormalTok{)}

\FunctionTok{lines}\NormalTok{(}
  \AttributeTok{x =}\NormalTok{ x,}
  \AttributeTok{y =} \FunctionTok{pdf\_ing\_gamma}\NormalTok{(x),}
  \AttributeTok{col =} \StringTok{"blue"}\NormalTok{,}
  \AttributeTok{lwd =} \DecValTok{2}
\NormalTok{)}

\FunctionTok{lines}\NormalTok{(}
  \AttributeTok{x =}\NormalTok{ x,}
  \AttributeTok{y =} \FunctionTok{pdf\_ing\_weibull}\NormalTok{(x),}
  \AttributeTok{col =} \StringTok{"darkgreen"}\NormalTok{,}
  \AttributeTok{lwd =} \DecValTok{3}
\NormalTok{)}

\FunctionTok{legend}\NormalTok{(}
  \StringTok{"topleft"}\NormalTok{,}
  \AttributeTok{legend =} \FunctionTok{c}\NormalTok{(}\StringTok{"Log{-}normal"}\NormalTok{, }\StringTok{"Gamma"}\NormalTok{, }\StringTok{"Weibull"}\NormalTok{),}
  \AttributeTok{col =} \FunctionTok{c}\NormalTok{(}\StringTok{"red"}\NormalTok{, }\StringTok{"blue"}\NormalTok{, }\StringTok{"darkgreen"}\NormalTok{),}
  \AttributeTok{lwd =} \DecValTok{2}
\NormalTok{)}
\end{Highlighting}
\end{Shaded}

\pandocbounded{\includegraphics[keepaspectratio]{04_distribution_fitting_files/figure-latex/unnamed-chunk-19-1.pdf}}

\subsubsection{Results}\label{results}

The histogram of the data was overlaid with theoretical curves
corresponding to each distribution.

\begin{Shaded}
\begin{Highlighting}[]
\NormalTok{x }\OtherTok{\textless{}{-}} \FunctionTok{seq}\NormalTok{(}
  \AttributeTok{from =} \FunctionTok{min}\NormalTok{(alior),}
  \AttributeTok{to =} \FunctionTok{max}\NormalTok{(alior),}
  \AttributeTok{length.out =} \DecValTok{1000}
\NormalTok{)}

\FunctionTok{hist}\NormalTok{(}
  \AttributeTok{x =}\NormalTok{ alior,}
  \AttributeTok{freq =} \ConstantTok{FALSE}\NormalTok{,}
  \AttributeTok{col =} \StringTok{"maroon"}\NormalTok{,}
  \AttributeTok{main =} \StringTok{"Histogram of Alior Bank Closing Stock Prices, 2023–2024"}\NormalTok{,}
  \AttributeTok{xlab =} \StringTok{"Closing price"}\NormalTok{,}
  \AttributeTok{ylab =} \StringTok{"Density"}\NormalTok{,}
  \AttributeTok{las =} \DecValTok{1}\NormalTok{,}
  \AttributeTok{breaks =} \DecValTok{30}
\NormalTok{)}

\FunctionTok{lines}\NormalTok{(}
  \AttributeTok{x =}\NormalTok{ x,}
  \AttributeTok{y =} \FunctionTok{pdf\_alior\_lnorm}\NormalTok{(x),}
  \AttributeTok{col =} \StringTok{"red"}\NormalTok{,}
  \AttributeTok{lwd =} \DecValTok{2}
\NormalTok{)}

\FunctionTok{lines}\NormalTok{(}
  \AttributeTok{x =}\NormalTok{ x,}
  \AttributeTok{y =} \FunctionTok{pdf\_alior\_gamma}\NormalTok{(x),}
  \AttributeTok{col =} \StringTok{"blue"}\NormalTok{,}
  \AttributeTok{lwd =} \DecValTok{2}
\NormalTok{)}

\FunctionTok{lines}\NormalTok{(}
  \AttributeTok{x =}\NormalTok{ x,}
  \AttributeTok{y =} \FunctionTok{pdf\_alior\_weibull}\NormalTok{(x),}
  \AttributeTok{col =} \StringTok{"darkgreen"}\NormalTok{,}
  \AttributeTok{lwd =} \DecValTok{2}
\NormalTok{)}

\FunctionTok{legend}\NormalTok{(}
  \StringTok{"topleft"}\NormalTok{,}
  \AttributeTok{legend=}\FunctionTok{c}\NormalTok{(}\StringTok{"Log{-}normal"}\NormalTok{, }\StringTok{"Gamma"}\NormalTok{, }\StringTok{"Weibull"}\NormalTok{),}
  \AttributeTok{col=}\FunctionTok{c}\NormalTok{(}\StringTok{"red"}\NormalTok{, }\StringTok{"blue"}\NormalTok{, }\StringTok{"darkgreen"}\NormalTok{),}
  \AttributeTok{lwd=}\DecValTok{2}
\NormalTok{)}
\end{Highlighting}
\end{Shaded}

\pandocbounded{\includegraphics[keepaspectratio]{04_distribution_fitting_files/figure-latex/unnamed-chunk-20-1.pdf}}

\begin{Shaded}
\begin{Highlighting}[]
\NormalTok{x }\OtherTok{\textless{}{-}} \FunctionTok{seq}\NormalTok{(}
  \AttributeTok{from =} \FunctionTok{min}\NormalTok{(ing),}
  \AttributeTok{to   =} \FunctionTok{max}\NormalTok{(ing),}
  \AttributeTok{length.out =} \DecValTok{1000}
\NormalTok{)}

\FunctionTok{hist}\NormalTok{(}
  \AttributeTok{x =}\NormalTok{ ing,}
  \AttributeTok{freq =} \ConstantTok{FALSE}\NormalTok{,}
  \AttributeTok{col =} \StringTok{"orange"}\NormalTok{,}
  \AttributeTok{main =} \StringTok{"Histogram of ING Bank Closing Stock Prices, 2023–2024"}\NormalTok{,}
  \AttributeTok{xlab =} \StringTok{"Closing price"}\NormalTok{,}
  \AttributeTok{ylab =} \StringTok{"Density"}\NormalTok{,}
  \AttributeTok{las =} \DecValTok{1}\NormalTok{,}
  \AttributeTok{breaks =} \DecValTok{30}
\NormalTok{)}

\FunctionTok{lines}\NormalTok{(}
  \AttributeTok{x =}\NormalTok{ x,}
  \AttributeTok{y =} \FunctionTok{pdf\_ing\_lnorm}\NormalTok{(x),}
  \AttributeTok{col =} \StringTok{"red"}\NormalTok{,}
  \AttributeTok{lwd =} \DecValTok{2}
\NormalTok{)}

\FunctionTok{lines}\NormalTok{(}
  \AttributeTok{x =}\NormalTok{ x,}
  \AttributeTok{y =} \FunctionTok{pdf\_ing\_gamma}\NormalTok{(x),}
  \AttributeTok{col =} \StringTok{"blue"}\NormalTok{,}
  \AttributeTok{lwd =} \DecValTok{2}
\NormalTok{)}

\FunctionTok{lines}\NormalTok{(}
  \AttributeTok{x =}\NormalTok{ x,}
  \AttributeTok{y =} \FunctionTok{pdf\_ing\_weibull}\NormalTok{(x),}
  \AttributeTok{col =} \StringTok{"darkgreen"}\NormalTok{,}
  \AttributeTok{lwd =} \DecValTok{2}
\NormalTok{)}

\FunctionTok{legend}\NormalTok{(}
  \StringTok{"topleft"}\NormalTok{,}
  \AttributeTok{legend =} \FunctionTok{c}\NormalTok{(}\StringTok{"Log{-}normal"}\NormalTok{, }\StringTok{"Gamma"}\NormalTok{, }\StringTok{"Weibull"}\NormalTok{),}
  \AttributeTok{col    =} \FunctionTok{c}\NormalTok{(}\StringTok{"red"}\NormalTok{, }\StringTok{"blue"}\NormalTok{, }\StringTok{"darkgreen"}\NormalTok{),}
  \AttributeTok{lwd    =} \DecValTok{2}
\NormalTok{)}
\end{Highlighting}
\end{Shaded}

\pandocbounded{\includegraphics[keepaspectratio]{04_distribution_fitting_files/figure-latex/unnamed-chunk-21-1.pdf}}

Preliminary visual analysis indicates that the Log-normal distribution
is the best-fitting model. Its curve follows the empirical frequency
points most closely, particularly in the central part of the
distribution.

However, all three models show significant deviations in the left tail.
The histogram indicates a higher concentration of observations in this
region, which none of the tested models fully capture. This suggests
that, while the overall fit is reasonable, the distributions may not
completely represent the characteristics of the data, especially for
extreme values.

\subsection{Goodness-of-Fit Analysis}\label{goodness-of-fit-analysis}

Interpretation of Goodness-of-Fit Statistics

To assess how well each theoretical distribution (Log-normal, Gamma,
Weibull) fits the empirical data, the following test statistics were
computed:

\begin{itemize}
\tightlist
\item
  \textbf{Kolmogorov--Smirnov (KS):} measures the maximum difference
  between the empirical and theoretical cumulative distribution
  functions (CDFs).\\
  Statistic:
\end{itemize}

\[
D_n = \sup_x |F_n(x) - F(x)|,
\]

where: - \(F_n(x)\) -- empirical cumulative distribution function
(ECDF),\\
- \(F(x)\) -- theoretical cumulative distribution function (CDF).

The smaller the value of \(D_n\), the better the fit.

\begin{itemize}
\tightlist
\item
  \textbf{Cramér--von Mises (CM):} integrates the squared difference
  between empirical and theoretical CDFs.

  \begin{itemize}
  \tightlist
  \item
    Lower values indicate a better fit.
  \end{itemize}

  Statistic:
\end{itemize}

\[
W^2 = n \int_{-\infty}^{\infty} (F_n(x) - F(x))^2 \, dF(x).
\]

\begin{itemize}
\tightlist
\item
  \textbf{Anderson--Darling (AD):} a weighted version of CM, giving more
  weight to the tails of the distribution.\\
  Statistic:
\end{itemize}

\[
A^2 = -n - \frac{1}{n} \sum_{i=1}^n \left( (2i-1)\big[\ln F(X_{(i)}) + \ln(1 - F(X_{(n+1-i)}))\big] \right),
\]

where \(X_{(i)}\) are the ordered observations.

The AD test places more emphasis on the fit in the tails of the
distribution compared to KS and CM.

\begin{itemize}
\item
  \textbf{Akaike's Information Criterion (AIC):} measures the trade-off
  between goodness of fit and model complexity.\\
  \[
  AIC = 2k - 2 \ln(\hat{L}),
  \]

  where:

  \begin{itemize}
  \tightlist
  \item
    \(k\) -- number of parameters,\\
  \item
    \(\hat{L}\) -- maximum value of the likelihood function.
  \end{itemize}

  Lower AIC → better fit.
\item
  \textbf{Bayesian Information Criterion (BIC):} similar to AIC but
  penalizes model complexity more strongly, especially with larger
  sample sizes.\\
  \[
  BIC = k \ln(n) - 2 \ln(\hat{L}),
  \]

  where:

  \begin{itemize}
  \tightlist
  \item
    \(n\) -- sample size.
  \end{itemize}

  Lower BIC → better fit.\\
  BIC penalizes models with more parameters more strongly than AIC.
\end{itemize}

\begin{Shaded}
\begin{Highlighting}[]
\FunctionTok{gofstat}\NormalTok{(}
  \AttributeTok{f =} \FunctionTok{list}\NormalTok{(alior\_fit\_lnorm\_rev, alior\_fit\_gamma\_rev, alior\_fit\_weibull\_rev),}
  \AttributeTok{fitnames =} \FunctionTok{c}\NormalTok{(}\StringTok{"Log{-}normal"}\NormalTok{, }\StringTok{"Gamma"}\NormalTok{, }\StringTok{"Weibull"}\NormalTok{),}
  \AttributeTok{discrete =} \ConstantTok{FALSE}
\NormalTok{)}
\end{Highlighting}
\end{Shaded}

\begin{verbatim}
## Goodness-of-fit statistics
##                              Log-normal      Gamma    Weibull
## Kolmogorov-Smirnov statistic  0.1225537  0.1329507  0.1316003
## Cramer-von Mises statistic    2.7962473  3.6375292  3.9571860
## Anderson-Darling statistic   18.7459398 22.0294383 23.5589019
## 
## Goodness-of-fit criteria
##                                Log-normal    Gamma  Weibull
## Akaike's Information Criterion   6424.787 6289.029 6281.107
## Bayesian Information Criterion   6433.976 6298.218 6290.296
\end{verbatim}

\begin{Shaded}
\begin{Highlighting}[]
\FunctionTok{gofstat}\NormalTok{(}
  \AttributeTok{f =} \FunctionTok{list}\NormalTok{(ing\_fit\_lnorm\_rev, ing\_fit\_gamma\_rev, ing\_fit\_weibull\_rev),}
  \AttributeTok{fitnames =} \FunctionTok{c}\NormalTok{(}\StringTok{"Log{-}normal"}\NormalTok{, }\StringTok{"Gamma"}\NormalTok{, }\StringTok{"Weibull"}\NormalTok{),}
  \AttributeTok{discrete =} \ConstantTok{FALSE}
\NormalTok{)}
\end{Highlighting}
\end{Shaded}

\begin{verbatim}
## Goodness-of-fit statistics
##                              Log-normal      Gamma   Weibull
## Kolmogorov-Smirnov statistic  0.1452921  0.1471288  0.149433
## Cramer-von Mises statistic    2.8299647  2.4322786  2.662052
## Anderson-Darling statistic   20.0217956 17.8290525 19.335855
## 
## Goodness-of-fit criteria
##                                Log-normal    Gamma  Weibull
## Akaike's Information Criterion   7997.883 7779.604 7749.787
## Bayesian Information Criterion   8007.072 7788.792 7758.976
\end{verbatim}

\subsubsection{Results and
Interpretation}\label{results-and-interpretation}

For both datasets, Alior Bank and ING, three theoretical distributions
were fitted: \textbf{Log-normal, Gamma, and Weibull}. Parameter
estimation was performed using \textbf{Maximum Likelihood Estimation
(MLE)}, and the quality of fit was assessed with \textbf{goodness-of-fit
tests} (Kolmogorov--Smirnov, Cramér--von Mises, Anderson--Darling) as
well as \textbf{information criteria} (AIC, BIC).

For \textbf{Alior Bank}, the goodness-of-fit tests indicate that the
\textbf{Log-normal distribution} provides the closest match to the
empirical data, particularly in the central part of the distribution. KS
measures the maximum deviation between the empirical and theoretical
cumulative distribution functions, CM averages the squared differences,
and AD gives more weight to the tails. Despite this, none of the
distributions fully capture the left tail, suggesting that extreme low
values are underrepresented. When considering \textbf{AIC and BIC},
which penalize model complexity while rewarding likelihood, the
\textbf{Weibull distribution} emerges as the most efficient overall
model. This highlights that different criteria can lead to different
conclusions: Log-normal fits the data better statistically, whereas
Weibull offers a better balance between fit and simplicity.

For \textbf{ING}, the goodness-of-fit tests favor the \textbf{Gamma
distribution}, which captures tail behavior more accurately. KS values
are similar across all distributions, while CM and AD highlight Gamma as
superior in describing the overall shape, including the tails. Again,
none of the models perfectly represent the left tail. However,
\textbf{AIC and BIC} select the \textbf{Weibull distribution} as the
most balanced model, reflecting an optimal compromise between fit and
parsimony.

\textbf{Overall}, no single distribution perfectly describes either
dataset. Goodness-of-fit tests emphasize closeness to empirical
distributions, especially in tails, while information criteria highlight
model efficiency and complexity. Combining these perspectives, one can
conclude that \textbf{Log-normal and Gamma distributions are better for
statistical fit}, whereas \textbf{Weibull provides the most balanced and
parsimonious model} across both datasets. Graphical inspection of
histograms and fitted curves reinforces these conclusions, particularly
regarding tail behavior.

\subsection{Hypothesis Testing of Distribution
Fit}\label{hypothesis-testing-of-distribution-fit}

The next step in the analysis is to formally test whether the fitted
distributions provide an adequate representation of the empirical data.
This is done by evaluating the \textbf{p-values} of the goodness-of-fit
tests. A p-value quantifies the probability of observing a test
statistic at least as extreme as the one computed from the sample,
assuming the null hypothesis is true.

For a \textbf{two-sided test}, the p-value is defined as:

\[
p\text{-value} = P(T \leq t_\text{obs}) + P(T \geq t_\text{obs})
\]

where \(T\) is the random variable representing the test statistic under
the null hypothesis, and \(t_\text{obs}\) is the observed value of the
statistic.

In this context, the \textbf{null hypothesis (\(H_0\))} states that the
empirical data come from the fitted distribution. If the p-value is
sufficiently high, we do not reject \(H_0\), which would suggest that
the distribution could plausibly describe the data.

However, there are important caveats for using p-values in the context
of \textbf{stock closing prices or log-returns}:

\begin{itemize}
\tightlist
\item
  Financial data often exhibit \textbf{heavy tails, skewness, and
  autocorrelation}, which standard goodness-of-fit tests may not fully
  account for.\\
\item
  A non-rejection of \(H_0\) (high p-value) does not guarantee that the
  model captures all important features of the data, particularly
  extreme events in the tails.\\
\item
  Conversely, very large sample sizes, as is common in financial time
  series, may produce extremely small p-values even for distributions
  that visually appear reasonable.
\end{itemize}

Thus, while p-values provide a formal test of fit, they should be
interpreted \textbf{together with graphical inspection and information
criteria}, rather than in isolation.

\begin{Shaded}
\begin{Highlighting}[]
\NormalTok{alior\_ks\_lnorm }\OtherTok{\textless{}{-}} \FunctionTok{ks.test}\NormalTok{(alior, }\StringTok{"plnorm"}\NormalTok{,}
                   \AttributeTok{meanlog =}\NormalTok{ alior\_fit\_lnorm\_rev}\SpecialCharTok{$}\NormalTok{estimate[}\StringTok{"meanlog"}\NormalTok{],}
                   \AttributeTok{sdlog   =}\NormalTok{ alior\_fit\_lnorm\_rev}\SpecialCharTok{$}\NormalTok{estimate[}\StringTok{"sdlog"}\NormalTok{])}
\end{Highlighting}
\end{Shaded}

\begin{verbatim}
## Warning in ks.test.default(alior, "plnorm", meanlog =
## alior_fit_lnorm_rev$estimate["meanlog"], : ties should not be present for the
## one-sample Kolmogorov-Smirnov test
\end{verbatim}

\begin{Shaded}
\begin{Highlighting}[]
\NormalTok{alior\_ks\_lnorm}
\end{Highlighting}
\end{Shaded}

\begin{verbatim}
## 
##  Asymptotic one-sample Kolmogorov-Smirnov test
## 
## data:  alior
## D = 0.68868, p-value < 2.2e-16
## alternative hypothesis: two-sided
\end{verbatim}

\begin{Shaded}
\begin{Highlighting}[]
\NormalTok{alior\_ks\_lnorm}\SpecialCharTok{$}\NormalTok{p.value }\SpecialCharTok{\textless{}} \FloatTok{0.05}
\end{Highlighting}
\end{Shaded}

\begin{verbatim}
## [1] TRUE
\end{verbatim}

\begin{Shaded}
\begin{Highlighting}[]
\NormalTok{alior\_ks\_lnorm}\SpecialCharTok{$}\NormalTok{p.value }\SpecialCharTok{\textless{}} \FloatTok{0.2}
\end{Highlighting}
\end{Shaded}

\begin{verbatim}
## [1] TRUE
\end{verbatim}

\begin{Shaded}
\begin{Highlighting}[]
\NormalTok{alior\_ks\_gamma }\OtherTok{\textless{}{-}} \FunctionTok{ks.test}\NormalTok{(}
\NormalTok{  alior, }\StringTok{"pgamma"}\NormalTok{,}
  \AttributeTok{shape =}\NormalTok{ alior\_fit\_gamma\_rev}\SpecialCharTok{$}\NormalTok{estimate[}\StringTok{"shape"}\NormalTok{],}
  \AttributeTok{rate  =}\NormalTok{ alior\_fit\_gamma\_rev}\SpecialCharTok{$}\NormalTok{estimate[}\StringTok{"rate"}\NormalTok{]}
\NormalTok{)}
\end{Highlighting}
\end{Shaded}

\begin{verbatim}
## Warning in ks.test.default(alior, "pgamma", shape =
## alior_fit_gamma_rev$estimate["shape"], : ties should not be present for the
## one-sample Kolmogorov-Smirnov test
\end{verbatim}

\begin{Shaded}
\begin{Highlighting}[]
\NormalTok{alior\_ks\_gamma}
\end{Highlighting}
\end{Shaded}

\begin{verbatim}
## 
##  Asymptotic one-sample Kolmogorov-Smirnov test
## 
## data:  alior
## D = 0.66956, p-value < 2.2e-16
## alternative hypothesis: two-sided
\end{verbatim}

\begin{Shaded}
\begin{Highlighting}[]
\NormalTok{alior\_ks\_gamma}\SpecialCharTok{$}\NormalTok{p.value }\SpecialCharTok{\textless{}} \FloatTok{0.05}
\end{Highlighting}
\end{Shaded}

\begin{verbatim}
## [1] TRUE
\end{verbatim}

\begin{Shaded}
\begin{Highlighting}[]
\NormalTok{alior\_ks\_gamma}\SpecialCharTok{$}\NormalTok{p.value }\SpecialCharTok{\textless{}} \FloatTok{0.2}
\end{Highlighting}
\end{Shaded}

\begin{verbatim}
## [1] TRUE
\end{verbatim}

\begin{Shaded}
\begin{Highlighting}[]
\NormalTok{alior\_ks\_weibull }\OtherTok{\textless{}{-}} \FunctionTok{ks.test}\NormalTok{(alior, }\StringTok{"pweibull"}\NormalTok{,}
                            \AttributeTok{shape =}\NormalTok{ alior\_fit\_weibull\_rev}\SpecialCharTok{$}\NormalTok{estimate[}\StringTok{"shape"}\NormalTok{],}
                            \AttributeTok{scale =}\NormalTok{ alior\_fit\_weibull\_rev}\SpecialCharTok{$}\NormalTok{estimate[}\StringTok{"scale"}\NormalTok{])}
\end{Highlighting}
\end{Shaded}

\begin{verbatim}
## Warning in ks.test.default(alior, "pweibull", shape =
## alior_fit_weibull_rev$estimate["shape"], : ties should not be present for the
## one-sample Kolmogorov-Smirnov test
\end{verbatim}

\begin{Shaded}
\begin{Highlighting}[]
\NormalTok{alior\_ks\_weibull}
\end{Highlighting}
\end{Shaded}

\begin{verbatim}
## 
##  Asymptotic one-sample Kolmogorov-Smirnov test
## 
## data:  alior
## D = 0.66099, p-value < 2.2e-16
## alternative hypothesis: two-sided
\end{verbatim}

\begin{Shaded}
\begin{Highlighting}[]
\NormalTok{alior\_ks\_weibull}\SpecialCharTok{$}\NormalTok{p.value }\SpecialCharTok{\textless{}} \FloatTok{0.05}
\end{Highlighting}
\end{Shaded}

\begin{verbatim}
## [1] TRUE
\end{verbatim}

\begin{Shaded}
\begin{Highlighting}[]
\NormalTok{alior\_ks\_weibull}\SpecialCharTok{$}\NormalTok{p.value }\SpecialCharTok{\textless{}} \FloatTok{0.2}
\end{Highlighting}
\end{Shaded}

\begin{verbatim}
## [1] TRUE
\end{verbatim}

\begin{Shaded}
\begin{Highlighting}[]
\NormalTok{ing\_ks\_lnorm }\OtherTok{\textless{}{-}} \FunctionTok{ks.test}\NormalTok{(ing, }\StringTok{"plnorm"}\NormalTok{,}
                        \AttributeTok{meanlog =}\NormalTok{ ing\_fit\_lnorm\_rev}\SpecialCharTok{$}\NormalTok{estimate[}\StringTok{"meanlog"}\NormalTok{],}
                        \AttributeTok{sdlog   =}\NormalTok{ ing\_fit\_lnorm\_rev}\SpecialCharTok{$}\NormalTok{estimate[}\StringTok{"sdlog"}\NormalTok{])}
\end{Highlighting}
\end{Shaded}

\begin{verbatim}
## Warning in ks.test.default(ing, "plnorm", meanlog =
## ing_fit_lnorm_rev$estimate["meanlog"], : ties should not be present for the
## one-sample Kolmogorov-Smirnov test
\end{verbatim}

\begin{Shaded}
\begin{Highlighting}[]
\NormalTok{ing\_ks\_lnorm}
\end{Highlighting}
\end{Shaded}

\begin{verbatim}
## 
##  Asymptotic one-sample Kolmogorov-Smirnov test
## 
## data:  ing
## D = 0.79373, p-value < 2.2e-16
## alternative hypothesis: two-sided
\end{verbatim}

\begin{Shaded}
\begin{Highlighting}[]
\NormalTok{ing\_ks\_lnorm}\SpecialCharTok{$}\NormalTok{p.value }\SpecialCharTok{\textless{}} \FloatTok{0.05} 
\end{Highlighting}
\end{Shaded}

\begin{verbatim}
## [1] TRUE
\end{verbatim}

\begin{Shaded}
\begin{Highlighting}[]
\NormalTok{ing\_ks\_lnorm}\SpecialCharTok{$}\NormalTok{p.value }\SpecialCharTok{\textless{}} \FloatTok{0.2}
\end{Highlighting}
\end{Shaded}

\begin{verbatim}
## [1] TRUE
\end{verbatim}

\begin{Shaded}
\begin{Highlighting}[]
\NormalTok{ing\_ks\_gamma }\OtherTok{\textless{}{-}} \FunctionTok{ks.test}\NormalTok{(ing, }\StringTok{"pgamma"}\NormalTok{,}
                        \AttributeTok{shape =}\NormalTok{ ing\_fit\_gamma\_rev}\SpecialCharTok{$}\NormalTok{estimate[}\StringTok{"shape"}\NormalTok{],}
                        \AttributeTok{rate  =}\NormalTok{ ing\_fit\_gamma\_rev}\SpecialCharTok{$}\NormalTok{estimate[}\StringTok{"rate"}\NormalTok{])}
\end{Highlighting}
\end{Shaded}

\begin{verbatim}
## Warning in ks.test.default(ing, "pgamma", shape =
## ing_fit_gamma_rev$estimate["shape"], : ties should not be present for the
## one-sample Kolmogorov-Smirnov test
\end{verbatim}

\begin{Shaded}
\begin{Highlighting}[]
\NormalTok{ing\_ks\_gamma}
\end{Highlighting}
\end{Shaded}

\begin{verbatim}
## 
##  Asymptotic one-sample Kolmogorov-Smirnov test
## 
## data:  ing
## D = 0.8239, p-value < 2.2e-16
## alternative hypothesis: two-sided
\end{verbatim}

\begin{Shaded}
\begin{Highlighting}[]
\NormalTok{ing\_ks\_gamma}\SpecialCharTok{$}\NormalTok{p.value }\SpecialCharTok{\textless{}} \FloatTok{0.05}
\end{Highlighting}
\end{Shaded}

\begin{verbatim}
## [1] TRUE
\end{verbatim}

\begin{Shaded}
\begin{Highlighting}[]
\NormalTok{ing\_ks\_gamma}\SpecialCharTok{$}\NormalTok{p.value }\SpecialCharTok{\textless{}} \FloatTok{0.2}
\end{Highlighting}
\end{Shaded}

\begin{verbatim}
## [1] TRUE
\end{verbatim}

\begin{Shaded}
\begin{Highlighting}[]
\NormalTok{ing\_ks\_weibull }\OtherTok{\textless{}{-}} \FunctionTok{ks.test}\NormalTok{(ing, }\StringTok{"pweibull"}\NormalTok{,}
                          \AttributeTok{shape =}\NormalTok{ ing\_fit\_weibull\_rev}\SpecialCharTok{$}\NormalTok{estimate[}\StringTok{"shape"}\NormalTok{],}
                          \AttributeTok{scale =}\NormalTok{ ing\_fit\_weibull\_rev}\SpecialCharTok{$}\NormalTok{estimate[}\StringTok{"scale"}\NormalTok{])}
\end{Highlighting}
\end{Shaded}

\begin{verbatim}
## Warning in ks.test.default(ing, "pweibull", shape =
## ing_fit_weibull_rev$estimate["shape"], : ties should not be present for the
## one-sample Kolmogorov-Smirnov test
\end{verbatim}

\begin{Shaded}
\begin{Highlighting}[]
\NormalTok{ing\_ks\_weibull}
\end{Highlighting}
\end{Shaded}

\begin{verbatim}
## 
##  Asymptotic one-sample Kolmogorov-Smirnov test
## 
## data:  ing
## D = 0.8283, p-value < 2.2e-16
## alternative hypothesis: two-sided
\end{verbatim}

\begin{Shaded}
\begin{Highlighting}[]
\NormalTok{ing\_ks\_weibull}\SpecialCharTok{$}\NormalTok{p.value }\SpecialCharTok{\textless{}} \FloatTok{0.05}
\end{Highlighting}
\end{Shaded}

\begin{verbatim}
## [1] TRUE
\end{verbatim}

\begin{Shaded}
\begin{Highlighting}[]
\NormalTok{ing\_ks\_weibull}\SpecialCharTok{$}\NormalTok{p.value }\SpecialCharTok{\textless{}} \FloatTok{0.2}
\end{Highlighting}
\end{Shaded}

\begin{verbatim}
## [1] TRUE
\end{verbatim}

\begin{longtable}[]{@{}llll@{}}
\toprule\noalign{}
Bank & Distribution & p-value \textless{} \(5\%\) & p-value \textless{}
\(20\%\) \\
\midrule\noalign{}
\endhead
\bottomrule\noalign{}
\endlastfoot
Alior Bank & Log-normal & TRUE & TRUE \\
Alior Bank & Gamma & TRUE & TRUE \\
Alior Bank & Weibull & TRUE & TRUE \\
ING & Log-normal & TRUE & TRUE \\
ING & Gamma & TRUE & TRUE \\
ING & Weibull & TRUE & TRUE \\
\end{longtable}

All goodness-of-fit tests performed for both \textbf{Alior Bank} and
\textbf{ING} yielded \textbf{p-values below conventional thresholds} (p
\textless{} 0.05 and p \textless{} 0.2). According to these tests, the
theoretical distributions differ significantly from the empirical data
in every case. In other words, \textbf{none of the fitted distributions
can be formally accepted as perfectly describing the observed data}.

Despite this, some models appear more plausible than others when
considering the relative performance across tests:

\begin{itemize}
\tightlist
\item
  For \textbf{Alior Bank}, the \textbf{Log-normal distribution} showed
  the closest alignment with the empirical data, both in central values
  and in the tails, compared to Gamma and Weibull.\\
\item
  For \textbf{ING}, the \textbf{Gamma distribution} was the most
  reasonable among the three, particularly in capturing tail behavior.
\end{itemize}

Visual inspection of the histograms suggests that \textbf{both datasets
might be better represented by a mixture of distributions}, reflecting
the presence of extreme values and possible multiple regimes. However,
modeling such mixed distributions would require a \textbf{much more
extensive analysis}, including identification of component distributions
and estimation of their parameters, which goes beyond the current scope.

\textbf{Conclusion:} While formal tests reject all simple distributional
models, Log-normal (Alior) and Gamma (ING) provide the most reasonable
approximations, serving as a practical baseline for further analysis,
albeit with clear limitations in the tails and extreme values.

\end{document}
